\chapter{Analytical Techniques for MAC Performance Evaluation}
\label{sec:chapter13}

This chapter compiles the mathematical and modeling frameworks used to evaluate wireless network performance at the MAC layer.

\section{Overview of Modeling Frameworks}
Contributions under the Mobile Computing Lab span several advanced mathematical domains:
\begin{itemize}
    \item \textbf{Markov Chains (Basker Deppuri, Sitaram):} Discrete-time and continuous-time Markov chains modeled after Bianchi's framework to analyze DCF throughput and collision probability.
    \item \textbf{Stochastic Methods (Arun, Hemanth):} Stochastic geometry and Poisson point processes to model spatial node distributions and interference profiles.
    \item \textbf{Fixed-Point Analysis (Aparna):} Solving non-linear systems of equations for multi-dimensional contention behavior.
    \item \textbf{Fluid Limit Approach (Aparna):} Asymptotic scaling and deterministic fluid limits of queue lengths.
    \item \textbf{Markov Decision Processes (MDP) (Senthil):} Sequential decision-making formulations for routing and power adaptation.
    \item \textbf{Queueing Theory (Anusree):} Modeling packet buffer dynamics, latency, and packet loss rate.
    \item \textbf{Age of Information (AoI) (Arun Kumar Reddy):} Characterizing fresh packet reception metrics and information update intervals in real-time systems.
\end{itemize}
